\section{第六条线：电子独有重联、能量去向与热化的边界}\label{story:energy}
磁通曲线和磁谱描述的只是重联的一部分。混合动理学研究显示，谱相近时离子参与方式仍可不同；电子独有重联既可关联小注入尺度，也可由大尺度重联产生的次级片或电子剪切生成。因此需要从“片有没有重联”继续追问“谁响应、谁受力、能量进了哪些自由度”。

这一线索连接三类容易混写的结果：Hall重排后的电流与电阻损失，电子平行相混合及弱碰撞热化，以及面内电场导致的离子加速和轨道展宽。共同的守恒约束是Hall项对总电流不直接做功；这不禁止它改变能量通道。以下分别核对分物种电功、内能和熵的证据，最后与下一节的实验原始诊断对接。

\subsection{电子独有重联：同样的能谱可隐藏不同的粒子参与}
Califano等的混合Vlasov计算以动理学离子和惯性等温电子比较两种湍流注入尺度。注入更靠近离子尺度的一组形成短电流片和电子独有喷流；另一组出现较长片以及离子、电子共同喷流。两组能谱却相似。这说明“谱看起来一致”不能单独验证流体模型的重联类型，必须同时检查离子响应、局地拓扑和片寿命。\citep{L267}

该文是\emph{动理学对照}而非纯Hall模拟。自动识别的数千电流峰并非全在重联，人工核对的重联层样本较小；电子惯性长度仅约1.36格，谱滤波与广义磁通的数值改变仍是限制。结果支持注入尺度可能选择电子独有重联的机制，不足以给出普适长度阈值。等温电子模型也不能用$\vct J\cdot\vct E'$直接预测电子升温。

\subsection{电子单独重联可以由大尺度重联和剪切流二次生成}
F655的2026年预印本\citep{F655}给出两条清楚的时序候选机制。其一，大尺度离子耦合重联先形成岛链，喷流使岛互相挤压，小于离子耦合尺度的新片再出现电子喷流。其二，湍流先生成电子速度剪切，随后出现与电子Kelvin--Helmholtz不稳定性相容的涡链，磁场被卷曲后形成小重联点。电子单独重联因此不必全靠直接在离子附近尺度注入能量。

这些图像有研究价值，但“太阳风中普遍发生”和“主导加热”尚超出本算例能证明的范围。系统仍为二维周期域，注入尺度约$50d_i$；电子尺度仅约一格。算法找到516个强电流结构，其中69\%薄于$d_i$且短于$5d_i$，并未逐个证实都是重联。部分$D_e$为负，所有片的平均量很小；不能把强局地正峰相加便当作净不可逆耗散。

结合解析速率研究\citep{F741}还得到另一个容易忽略的现象：无离子喷流并不意味着离子不反馈。周期小系统的离子反压仍可通过面内电场限制电子平台速度，使远上游离子归一率为$O(1)$；电子峰值、平台速度、局地电子归一率是三个不同量。

\subsection{回旋平均、能谱与加热：全局相似不意味着模型等价}
\citet{E471}的四场陀螺流体显示，热离子改变导心流、X点尖角、电子密度分离线结构和高谐波，而部分全局磁通仍相近。
其$N_i$是导心密度，不是仪器测到的粒子密度；把两者直接比较会把极化响应误判成电子--离子非准中性。
\citet{E474}加入有限$\beta_e$和电子FLR后出现非线性减速间隔与第二轮超指数阶段，电子密度出现嵌套四极及更平缓梯度，而广义不变量仍继续细丝化。
此结果的方向随固定参数组改变，见式\eqref{eq:selected-beta}；部分增长率振荡还来自X点移动。

\citet{L264}比较Hall与混合动理学湍流时，磁场谱相似而速度及能量分配可不同；对比时段对应不同成熟阶段，不能仅由频率相近建立完整因果能量通路。
\citet{F644}的KH驱动Hall谱在离子尺度附近变陡，并在小电流片识别重联候选；但模型没有显式电子破冻结项，数值滤波就是必须审计的部分。
这些研究支持Hall模型捕获部分多尺度磁结构，尚未证明它精确预测粒子热化。

\citet{E512}提供了一个清楚的边界：电子平行相混合可把流体自由能送到高Hermite矩，再由弱碰撞热化。
这需要流体--动理学扩展，不能用纯Hamiltonian双流体的磁能下降代替。
\citet{E553}又从分布重建说明有限矩的内容边界：即使21矩恢复三角形主体，仍缺双峰与精细多轨道群。
\citet{L313}所讨论的电子正则磁场冻结则强调惯性拓扑变化不必耗散能量；它是理论框架讨论，没有本篇新的大规模重联模拟。
三维指数分离可缩短特定流程中的尺度形成时间，但三维并不自动保证混沌，常密度固定离子前提也不可随意推广。

\subsection{“Hall电流耗散”不是Hall项直接加热}
该文发现面内电流的电阻耗散$\eta(J_x^2+J_y^2)$占总$\eta J^2$的比例可从约40\%升到55\%，沿场电流的比例约45\%。这两种划分并不相互补足，因为面内／面外与平行／垂直是不同坐标分解。Hall项自身依然满足$\J\cdot(\J\times\B)=0$。变化的是Hall重排后的电流结构及其\emph{电阻}损失，不能称Hall项把一半磁能直接变成热。

原方程压强绝热演化，没有把$\eta J^2$回收到内能，因此这些量也不是完整的分物种加热预测。中心线磁谱从浅段加陡段过渡到近单一约$k^{-2.22}$，但固定线谱、有限拟合区间和不同平均时间窗还不足以证明普适湍流级联律。\citep{L263}

\subsection{Hall电场与离子加热：轨道证据应怎样读}\label{sec:synthesis-heating}
\citet{L272}把电子正则涡量描述、给定场试粒子与PIC图像连接起来，提出强面内电场梯度使离子失去规则回旋运动的机制。所示PIC中，梯度指标较大的区域与高离子温度区大致对应；导引场由0增至0.5时，强效应的范围和程度减小。它提供了比“释放了磁能所以会热”更具体的空间预测，但图6、7来自其2019前序工作，不能计为一套独立复现。

认真读方法还需保留三个区别。第一，试粒子图1、2把粒子\emph{能量分布的标准差}记作TE，它不是通常由中心速度二阶矩定义的温度。第二，倒抛物势的线性运动方程本来就有指数发散解；仅凭这个例子不能独立建立复杂混沌或不可逆熵增。第三，论文中由准中性采用不可压缩，以及把给定几何下的阈值称作普遍热化条件，都需要额外假设。第一章第\ref{sec:synthesis-heating-model}节给出带符号的轨道方程、倒抛物势系数复核和能量分解。

这些限制没有取消其PIC中的空间关联，却改变了结论强度：它支持一条可检验的失磁化与增能途径，尚不足以单独解释所有导场、三维和高$\beta$重联中的热化。与\citet{E512}的相混合—弱碰撞链相比，有限矩模型还缺乏对速度空间细化及熵产的直接约束；与第\ref{sec:synthesis-mrx}节MRX原文比较，也不能把后续解释倒写为实验当时已唯一证实的机制。

\subsection{双温Hall已经能算局地加热，分配可信度仍取决于闭合}
2026年11月12日的加热专题沿上述问题继续核读。\citet{L223}图3给出MAST双温Hall计算：离子黏性耗散在下游形成双峰，电子耗散空间分布不同；其前身Stanier2013采用固定$p_i=p_e=p/2$，不能把前身的结果当成独立分配预测。L223双温扩展仍是初步结果，电子热点受自由超电阻参数控制，图3网格/时间步与独立收敛序列未报告。\citet{Kaptanoglu2020}进一步以三维双温Hall研究HIT-SI驱动、温度与损失并比较多种实验诊断；存在尚未取得全文的勘误，不能把未经勘误对照的精确比例作为核心结论。

这两篇直接说明“Hall模型完全不能研究离子/电子加热”过强。可以计算的是\emph{指定闭合下}的局地温度和能量预算；能否代表真实热化，须检验黏性、热流、物种交换及电阻/超电阻。\citet{L164}用共同流速、独立两温的碰撞激波模型得到离子先升温、电子后经碰撞和热传导响应；它不是自洽求Hall重联层。\citet{L247}的一般双流体方程含完整输运，但第IV节四个验证算例强制两温相等并删去离子能量方程，不能将其数值验证说成加热比验证。

\subsection{从轨道展宽到两物种熵产：补全前后文仍不能混算模型}
\citet{Yoon2018Heating}先以固定离子EMHD产生场，再积分测试离子；\citet{F743}用自洽PIC支持低/中导场下的失磁化机制。L272的图6--7继承后者，不能算独立再验证。\citet{Yoon2023PhaseSpace}进一步把冷多束叠加导致的高二阶矩与混沌展宽拆开；作者仍明确留下磁能进入离子的精确份额，以及有序/随机增能相对份额。该链没有演化出普适Hall两温闭合。

另一方面，\citet{Loureiro2013Heating}基于E512计算电子相混的高阶Hermite自由能与延迟耗散；它没有完整离子竞争通道。\citet{Numata2015Heating}用双物种陀螺动理学给出随$\beta$增加离子熵产作用增强的条件性结果，部分晚期瞬时$D_i/D_e$达量级1；这既不是累计普适比例，也不是有限矩Hall的成绩。L244的有限矩Landau模型实际用冷离子，耗散峰没有上述滞后；这正是闭合需要检验的内容。\citep{L244}

\citet{TenBarge2019MMS}图1已直接输出两物种流能和热能，十矩二维/三维全域预算相近；但热能增量没有扣除所有可逆压缩。电子各向异性幅度接近MMS、平行优先升温区却落在不同侧，固定尺度闭合又抑制LHDI。\citet{Finelli2021}证明改善流体电子闭合会反馈动理学离子，但作者把详细能量分配留待未来；\citet{Che2014Buneman}是PIC诊断后的近似双流体描述，实际未发生重联，仍需动理学理论或PIC标定拖曳。模型身份不能因标题含“two-fluid”而合并。

\subsection{实验的2:1究竟是哪一个比例}
\citet{F679}扩展Yoo2013，增加实验量级离子碰撞VPIC：无碰撞下游包含多束，加入碰撞后更近Maxwell。它支持电势加速、轨道混合到下游热化的条件性路径。\citet{F674}正式版图7中离子/电子沉积功率为$0.71\pm0.15$/$0.35\pm0.06$ MW，约2:1；同图的热能变化却为$0.27\pm0.10$/$0.26\pm0.10$ MW，约1:1，而且按式(6)包含储存与边界焓流，仍不是局地熵产。\citet{Yoo2017Heating}非对称MRX两种沉积更相近，直接限制普适2:1的外推。

\citet{Tanabe2015MAST}测到窄电子热点、下游离子热点及约4--10 ms物种交换时间，直接引用L223作温度结构解释；缺乏直接磁探针约束，不能唯一确定黏性/激波机制。强导场、闭合磁面与驻留时间也不容直接套用弱导场PIC趋势。\citet{TS6FEC2025}虽提及Hall MHD，实际报告的模拟是VPIC，约20\%比较针对电势幅度；独立Hall加热结果未展示，不能补作加热比验证。

\subsection{拾取与轨道混合：补上的参照与未统一的归因}
\citet{Drake2009Pickup}（2009，JGR）把离子跨窄Hall磁场边界后的非绝热运动、反向束和下游温度联系起来。图2给轨道与分布，图3在随出流的参考系中消去主要电场；冷入流近似下式(8)给$\Delta T_i\simeq m_iv_0^2/3$。该温度是中心二阶矩，不要求Maxwell分布。图6的22次Wind高剪切出流观测支持速度平方标度，但系数仅预测的0.39；各向异性也未定量吻合。其正式版14页含附录A已在11月13日新增核读，不能与尚未核读的2009年ApJ强导场/丰度论文混为一篇。

Drake2009第16段在下游近稳态轨道模型中认为面内静电场往返功相消，净增能与重联电场有关；\citet{Yoon2023PhaseSpace}原文Discussion部分和结论却批评以拾取解释零/弱导场整体增能，强调单阱混合与面内场，并称重联电场净功为零。这是原文之间尚未由本报告裁决的差异。其图7分力的速度积分不能单独代替逐分量$\int q_sE_av_a\dd t$的能量积分，末速度恢复亦不单独证明累计电功为零。应在同一场、粒子和时间窗核对净动能与全部电功；不能将这两种叙述直接写成统一的Hall热化定律。

\paragraph{这一线索走到了哪里。}双温Hall/双流体已经能研究条件性的两物种温度和能量预算；电子独有喷流、磁能下降、正局地电功与不可逆热化仍是不同量。应分别记录温度比、沉积率比、内能变化和熵产，核查同一控制体的压力功、热流与边界损失。独立专题补充场/流几何、轨道与相混物理图像，展开诊断和逐图阅读，并将已修复的报告遗漏与原文未决分别登记。文件位于项目根目录\texttt{specific problems/离子加热问题/}。
